% ══════════════════════════════════════════════════════════════════════════════
\chapter{Introduction}
\label{ch:introduction}
% ══════════════════════════════════════════════════════════════════════════════

\begin{figure}[H]
    \centering
    \includegraphics[width=0.85\textwidth]{photos/station_wide.jpg}
    \caption{Assembled, field-deployed ground station tracking a pass.}
\end{figure}

\paragraph{Context.}
A CubeSat is a small standardised satellite built up from 10~cm cubic units~\cite{cubesat-spec}; a 3U platform fits inside a 30$\times$10$\times$10~cm envelope, weighs a few kilograms, and can carry a meaningful scientific payload to low Earth orbit (LEO) on a rideshare launch for a fraction of the cost of a dedicated mission. KU~Leuven's AetherSpace programme is one such project, targeting a heat-shield sample-return demonstration with launch around 2030. The flight transceiver is a Texas Instruments CC1200~\cite{cc1200} on a 433~MHz UHF downlink at +14~dBm ($\approx$25~mW) into a 2~dBi turnstile or dipole antenna. Once in orbit a CubeSat is only as useful as the ground infrastructure that receives its telemetry, and providing that ground segment for AetherSpace is the subject of this thesis.

\paragraph{SatNOGS.}
The natural ecosystem to plug into is SatNOGS (Satellite Networked Open Ground Station), an open-source network of amateur ground stations that share scheduling and decoded telemetry through a public web platform. Each participating station tracks the satellites it can see, decodes the frames it captures, and uploads them to a central database; AetherSpace mission operations can then query that database without owning the receiving hardware. A station joins the network by speaking two standard interfaces: EasyComm~III over serial to the rotator and SiDS HTTP to the telemetry database. Beyond those two contracts, the choice of how to build the station is left to the operator. A full ecosystem walk-through, including the SatNOGS Network scheduler, the SDR-based receive client, and the comparison against earlier centralised projects such as GENSO, is given in the 2023 KU~Leuven masterproef by Hanssens~\cite{hanssens2023}, on which this work builds; \S\ref{sec:satnogs} summarises the parts of the ecosystem that are directly relevant to the design decisions in this thesis.

\paragraph{The gap.}
University CubeSat programmes are proliferating, but the ground-station hardware that should accompany them is not. Commercial AZ/EL rotators able to track a 433~MHz LEO link cost \euro{}760--1300~\cite{yaesu-g5500,rfham-spid}; the open-source SatNOGS v3 reference design itself still requires \euro{}300--500 in parts~\cite{satnogs-v3,satnogs-v3-bom}, and assumes a permanent rooftop installation that a student team without dedicated facilities cannot easily realise. For a project the size of AetherSpace, a fixed installation is overkill and the budget cannot justify a commercial unit.

\paragraph{This thesis.}
This work designs, builds, and field-validates a portable SatNOGS-compatible ground station that targets approximately \euro{}100 in station-core components, deploys from a standard camera tripod in about four minutes, and integrates with the SatNOGS rotator ecosystem through the standard EasyComm~III/hamlib interface. The target application is the AetherSpace mission, but neither the rotator nor the Pi-side software stack is AetherSpace-specific: any satellite present in the SatNOGS DB can be tracked, and any operator with a hamlib-compatible rotator can reuse the software.

\section{Problem Statement}
\label{sec:problem-statement}

The chapter opener established the mission (AetherSpace), the ecosystem (SatNOGS), and the gap (no cheap, portable, student-buildable station exists). This section turns those into engineering constraints. Five aspects of the problem set the design envelope: the geometry of a LEO pass, the link budget against the AetherSpace transmitter, the cost ceiling, the AetherSpace target itself, and the radio-architecture choice that follows from picking the same transceiver at both ends.

\paragraph{Geometry.}
A LEO satellite at 400~km altitude is visible from a given ground station for 5--8~minutes per 92-minute orbit. The total angular rate at zenith is bounded by
\begin{equation}
    \dot{\theta}_{\max} = \frac{v}{h} = \frac{7.67~\text{km/s}}{400~\text{km}} = 0.01918~\text{rad/s} = 1.10~\text{deg/s},
\end{equation}
with $v = \sqrt{\mu/(R_E+h)} = 7.67$~km/s for $\mu = 3.986 \times 10^{14}$~m$^3$/s$^2$ and $R_E+h = 6771$~km~\cite{montenbruck2000orbits}. For passes with closest-approach elevation above $\sim$70\degree\ the instantaneous azimuth rate exceeds 10\degree/s for a few seconds, because a small displacement of the sub-satellite point maps to a large arc in azimuth.

\paragraph{Link.}
The Siretta Oscar 44 Yagi used in this work (9~dBi peak, 54\degree$\times$48\degree\ 3-dB beamwidths at 434~MHz~\cite{siretta-oscar44}) is required to close the link against 145.8~dB of free-space path loss at 433~MHz and 30\degree\ elevation (Chapter~\ref{ch:background}, Table~\ref{tab:link-budget}). Its narrow beam sets the pointing-accuracy requirement on the rotator~\cite{stutzman2012antenna}.

\paragraph{Cost.}
Commercial rotators cost \euro{}760--900 (Yaesu G-5500~\cite{yaesu-g5500}) to \euro{}1200--1300 (SPID RAS~\cite{rfham-spid}); the SatNOGS v3 open-source rotator costs \euro{}300--500 in parts~\cite{satnogs-v3, satnogs-v3-bom}. The v2 MOTOR\_RF\_HAT BOM derived here (\S\ref{sec:cost}, design-complete) is approximately \euro{}40 for the combined controller/RF PCB, \euro{}75 for the full rotator hardware (PCB + motors + mechanics, excluding the Pi), and \euro{}100 for the station core (excluding antenna, coax, power supply, tripod).

\paragraph{Target.}
The full AetherSpace specification is given in \S\ref{sec:aetherspace}. The constraint it imposes here is link-budget closure: against a +14~dBm transmitter into a 2~dBi turnstile, the link in Chapter~\ref{ch:background} closes with positive theoretical margin only near zenith, and at 30\degree\ elevation the theoretical margin is $-5.7$~dB (Table~\ref{tab:link-budget}). Mid-elevation passes therefore demand both a directional receive antenna and a low-noise front end.

\paragraph{Radio choice.}
The station uses the CC1200 transceiver~\cite{cc1200} at both ends. This deviates from the canonical SatNOGS receive chain (SDR + GNU Radio software demodulation) but eliminates per-satellite modulation-parameter uncertainty: both ends share the same SmartRF register set once the flight configuration is finalised.

\subsection*{A deliberate alternative, not a replacement}

The canonical SatNOGS receive chain is a Raspberry Pi running the official \code{satnogs-client} daemon, an RTL-SDR v3 dongle feeding GNU Radio flowgraphs (\code{gr-satnogs}) through SoapySDR, and a hamlib-compatible rotator~\cite{satnogs-main, hamlib}. That stack decodes a broad range of amateur-satellite modulations (FSK, GFSK, AFSK, BPSK, LoRa chirp, CW, SSTV) in software, is maintained by the Libre Space Foundation, and remains the more flexible choice for a general-purpose amateur-radio ground station. The rotator produced by this thesis plugs into that stack unchanged over EasyComm III (hamlib model 204~\cite{jackson1990easycomm}).

This thesis adopts the CC1200 hardware-packet-modem path for three reasons.

The first is mission alignment. AetherSpace carries a CC1200 on-board, and using the same transceiver at the ground end eliminates modulation-parameter matching uncertainty: a single SmartRF Studio register set loaded at both ends configures identical symbol rate, deviation, filter bandwidth, sync word, and packet framing. An SDR receiver would need those values re-measured or re-derived per satellite and per flight revision.

The second is educational scope. The thesis is a joint MSc in electronics, and the CC1200 path exercises a full embedded-systems stack end-to-end: PCB design with controlled-impedance RF routing and matching networks (Chapter~\ref{ch:rf-frontend}), RP2040 firmware with SPI drivers and a binary COBS protocol (\S\ref{sec:rf-hat-firmware}, with protocol structure in Appendix~\ref{app:protocol}), Pi-side Python controlling CC1200 registers over UART (\S\ref{sec:pi-stack}), and a closed-loop tracking rotator (Chapters~\ref{ch:system-design} and~\ref{ch:firmware-software}).

The third is form-factor fit. A Pi 3A+ ships with 512~MB of LPDDR2 SDRAM~\cite{rpi3aplus}, of which approximately 416~MB is available to userspace after the default 64~MB GPU split. The SatNOGS Docker stack targets at least 1~GB of RAM~\cite{satnogs-main}; the Python stack described in \S\ref{sec:pi-stack} ($\sim$7900 LOC across the Pi packages listed in Appendix~\ref{app:bom}) runs on the 3A+ within a 200--300~MB resident-set footprint measured with \code{top}, leaving margin for the HTTPS dashboard and the systemd services.

\section{Research Questions}

The central question of this thesis is:

\begin{quote}
    Can an economical, portable SatNOGS-compatible ground station be designed, built, and field-validated as the ground segment for the AetherSpace CubeSat mission?
\end{quote}

This breaks down into six sub-questions:

\begin{enumerate}
    \item Can a satellite-tracking rotator be built for under \euro{}50 in materials?
    \item What pointing accuracy is achievable with geared DC motors, 3D-printed structure, and encoder-based PID control?
    \item Can the system integrate with the SatNOGS network using only the standard EasyComm/hamlib interface?
    \item What are the empirical trade-offs of DC motors versus stepper motors for satellite tracking?
    \item What sensitivity does the assembled CC1200 receive chain achieve in practice, and what is the dominant factor limiting it relative to the datasheet specification?
    \item Is a CC1200 hardware-modem approach a viable alternative to the SDR-based SatNOGS receive chain for the AetherSpace mission?
\end{enumerate}

\section{Scope and Division of Work}

\begin{itemize}
    \item Stan Coene: Mechanical rotator design, motor controller PCB (RP2040 + dual motor driver), rotator firmware, Raspberry Pi software stack, systemd services, web dashboard, field deployment workflow, system integration.
    \item Robbe Lehaen: RF front-end hardware (dual CC1200 transceiver HAT, UHF 433~MHz assembled + VHF 144/145~MHz routed), matching network design, Assembly of PCB's, RF HAT firmware (COBS protocol, SPI driver), CC1200 register configuration.
\end{itemize}

\section{Thesis Outline}

\begin{description}[style=unboxed, leftmargin=0pt]
    \item[Chapter~\ref{ch:background}] derives the physics of LEO satellite communication from first principles: orbital mechanics, link budgets, Doppler shift, and the SatNOGS ecosystem.
    \item[Chapter~\ref{ch:system-design}] presents the system architecture, mechanical design, rotator controller PCB, power distribution, bring-up lessons, and cost analysis.
    \item[Chapter~\ref{ch:rf-frontend}] describes the dual-CC1200 RF front-end HAT: PCB design, layout, and matching networks.
    \item[Chapter~\ref{ch:firmware-software}] covers the rotator firmware, RF HAT firmware, Raspberry Pi software stack, and field-deployment workflow.
    \item[Chapter~\ref{ch:results}] gives results and validation: tracking accuracy, RF measurements, field tests, and reception analysis. The link-budget derivation itself lives in Chapter~\ref{ch:background} and is only referenced here.
    \item[Chapter~\ref{ch:conclusion}] provides conclusions, limitations, and future work.
\end{description}
